\section*{Capítulo \ref{ch:g10:habitable-planet} --- Un planeta habitable}

\begin{solution}{exo:g10:habitable-planet:1}
Es la temperatura a la que un emisor perfecto radiaría hacia el espacio
exactamente la potencia que el planeta absorbe de su estrella. La superficie
de la Tierra está más caliente (\qty{288}{K} frente a unos \qty{255}{K})
porque los gases de efecto invernadero absorben parte del infrarrojo que
emite y devuelven una parte hacia abajo.
\end{solution}

\begin{solution}{exo:g10:habitable-planet:2}
Venus: \qty{737}{K} medidos frente a una \omterm{def:g10:habitable-planet:equilibrium-temperature}{temperatura de equilibrio} de unos
\qty{227}{K}, un calentamiento de unos \qty{510}{K}.
\end{solution}

\begin{solution}{exo:g10:habitable-planet:3}
$S = 1361/1.52^2 \approx \qty{589}{W/m^2}$ (la ficha de datos da
\qty{586}{W/m^2}, con la distancia media exacta).
\end{solution}

\begin{solution}{exo:g10:habitable-planet:4}
\qty{1}{bar}, \qty{330}{K}: líquida (por debajo del punto de ebullición,
unos \qty{373}{K}). \qty{0.001}{bar}, \qty{250}{K}: hielo (la presión queda
por encima de la curva hielo--vapor, que está a unos \qty{0.0008}{bar} a
\qty{250}{K}). \qty{100}{bar}, \qty{700}{K}: vapor; la temperatura supera la
temperatura crítica, así que no puede existir líquido.
\end{solution}

\begin{solution}{exo:g10:habitable-planet:5}
Un planeta que gira alrededor de una estrella distinta del Sol. Tránsitos:
el planeta pasa por delante de su estrella y la atenúa ligeramente en cada
órbita. Oscilación de la estrella: la atracción del planeta la desplaza
adelante y atrás a lo largo de nuestra visual, lo que desplaza
periódicamente el color de su luz.
\end{solution}

\begin{solution}{exo:g10:habitable-planet:6}
$T = (1361 \times 0.706/(4 \times 5.67\times10^{-8}))^{1/4} \approx
\qty{255}{K} \approx \qty{-18}{\degreeCelsius}$. Con $A = 0$:
$\approx \qty{278}{K} \approx \qty{5}{\degreeCelsius}$.
\end{solution}

\begin{solution}{exo:g10:habitable-planet:7}
$\sqrt 4 = 2$: de $2 \times 0.99 = \qty{1.98}{au}$ a $2 \times 1.70 =
\qty{3.40}{au}$. A \qty{1}{au}, el planeta recibiría cuatro veces la
irradiancia de la Tierra: muy por dentro del borde interior, demasiado
caliente.
\end{solution}

\begin{solution}{exo:g10:habitable-planet:8}
$S = 1361/0.723^2 \approx \qty{2604}{W/m^2}$;
$T_{\mathrm{eq}} = (2604 \times 0.23/(4\sigma))^{1/4} \approx \qty{227}{K}$.
La superficie está unos \qty{510}{K} más caliente: su densa atmósfera de
dióxido de carbono (\qty{92}{bar}) produce un efecto invernadero enorme.
\end{solution}

\begin{solution}{exo:g10:habitable-planet:9}
A \qty{6.36}{mbar}, el agua solo es líquida entre unos
\qty{0}{\degreeCelsius} y \qty{0.5}{\degreeCelsius}. A
\qty{10}{\degreeCelsius}, la presión queda por debajo de la curva de
ebullición: el agua se evapora hirviendo; de noche se congela.
\end{solution}

\begin{solution}{exo:g10:habitable-planet:10}
$11.2/2.4 \approx 4.7$. Una molécula de gas tiene que moverse casi cinco
veces más deprisa para escapar de la Tierra que de la Luna; a lo largo de
miles de millones de años, la gravedad débil de la Luna dejó escapar sus
gases, mientras que la de la Tierra retuvo la mayor parte.
\end{solution}

\begin{solution}{exo:g10:habitable-planet:11}
$T = (1361 \times 0.50/(4\sigma))^{1/4} \approx \qty{234}{K}$, unos
\qty{-39}{\degreeCelsius}. El hielo reflejaría tanta luz solar que el
planeta seguiría lo bastante frío para conservarlo: un bucle que se refuerza
a sí mismo.
\end{solution}

\begin{solution}{exo:g10:habitable-planet:12}
$T_{\mathrm{eq}} = (S_0(1-A)/(4\sigma d^2))^{1/4} = (S_0/(4\sigma))^{1/4}
(1-A)^{1/4} (d^2)^{-1/4}$; $(1361/(4\sigma))^{1/4} \approx \qty{278.3}{K}$ y
$(d^2)^{1/4} = \sqrt d$. Cuatro veces más lejos: $T$ se divide entre $\sqrt 4 = 2$.
\end{solution}

\begin{solution}{exo:g10:habitable-planet:13}
$\sqrt d = 278.3 \times 0.706^{1/4}/273 \approx 0.934$, luego $d \approx
\qty{0.873}{au}$, por dentro del borde interior (\qty{0.99}{au}). La
\omterm{def:g10:habitable-planet:equilibrium-temperature}{temperatura de equilibrio} no tiene en cuenta el efecto invernadero: un
planeta real con atmósfera está más caliente que $T_{\mathrm{eq}}$, así que
se sobrecalienta y pierde su agua ya más lejos, hacia \qty{0.99}{au}. Los
bordes salen de modelos de la atmósfera completa, no de $T_{\mathrm{eq}}$
sola.
\end{solution}

\begin{solution}{exo:g10:habitable-planet:14}
Entre el punto de fusión, unos \qty{0}{\degreeCelsius}, y el punto de
ebullición, unos \qty{7}{\degreeCelsius}. A \qty{15}{\degreeCelsius}, el
lago herviría.
\end{solution}

\begin{solution}{exo:g10:habitable-planet:15}
$\sqrt{0.5} \approx 0.707$: de $0.707 \times 0.99 \approx \qty{0.70}{au}$
a $0.707 \times 1.70 \approx \qty{1.20}{au}$. El planeta está justo en el
borde exterior. Los bordes son resultados de modelos y el escalado es
aproximado, así que es imposible decir si está dentro o fuera.
\end{solution}

\begin{solution}{pb:g10:habitable-planet:1}
\textbf{1.} $0.00151/0.0485^2 \approx 0.642$.

\textbf{2.} $0.642 \times 1361 \approx \qty{874}{W/m^2}$.

\textbf{3.} Entre ambas: unas tres veces menos que Venus y una vez y media
más que Marte.

\textbf{4.} $365.25/11.18 \approx 32.7$.

\textbf{5.} $(874/(4 \times 5.67\times10^{-8}))^{1/4} \approx \qty{249}{K}$.

\textbf{6.} $(874 \times 0.23/(4\sigma))^{1/4} \approx \qty{173}{K}$.

\textbf{7.} Ninguna: \qty{249}{K} sigue por debajo de \qty{273}{K}. Una
cubierta brillante de nubes o de hielo baja la \omterm{def:g10:habitable-planet:equilibrium-temperature}{temperatura de equilibrio} en
más de \qty{75}{K}: lo que cubre un planeta importa tanto como su distancia.

\textbf{8.} $\sqrt{0.00151} \approx 0.0389$.

\textbf{9.} $0.99 \times 0.0389 \approx \qty{0.0385}{au}$.

\textbf{10.} $1.70 \times 0.0389 \approx \qty{0.066}{au}$.

\textbf{11.} Sí: $0.0385 < 0.0485 < 0.066$.

\textbf{12.} $d_\odot = 0.0485/0.0389 \approx \qty{1.25}{au}$: en la mitad
interior de la \omterm{def:g10:habitable-planet:habitable-zone}{zona habitable} del Sol, entre la Tierra y Marte.

\textbf{13.} $1.055^{1/3} \approx 1.018$.

\textbf{14.} $11.2 \times \sqrt{1.055/1.018} \approx \qty{11.4}{km/s}$,
más o menos la de la Tierra: su gravedad podría retener una atmósfera (que
la tenga o no depende de su historia).

\textbf{15.} La oscilación de la estrella solo revela la parte de su
movimiento a lo largo de nuestra visual; como no se conoce la inclinación de
la órbita, la masa obtenida solo es la verdadera si la órbita se ve de
canto, y en los demás casos es un mínimo.

\textbf{16.} La presión del punto triple, \qty{611.657}{Pa}, unos
\qty{0.006}{bar}.

\textbf{17.} $(874 \times 0.706/(4\sigma))^{1/4} \approx \qty{228}{K}$.

\textbf{18.} $288 - 228 = \qty{60}{K}$: Próxima~b necesitaría un
calentamiento por efecto invernadero de \textbf{unos \qty{60}{K}}, casi el
doble de los \qty{33}{K} de la Tierra.
\end{solution}
